\par\noindent\textbf{(1)} 取K₀=|0⟩⟨0|+√η|1⟩⟨1|、K₁=√(1−η)|0⟩⟨1|。直接算 K₀†K₀+K₁†K₁=|0⟩⟨0|+[η+(1−η)]|1⟩⟨1|=I，因此映射 E(ρ)=Σ\_iK\_iρK\_i†迹保持；Kraus形式自动保证完全正性。
\par\noindent\textbf{(2)} 令ρ=[ρ₀₀,ρ₀₁;ρ₁₀,ρ₁₁]，代入得 E(ρ)=[ρ₀₀+(1−η)ρ₁₁,√ηρ₀₁;√ηρ₁₀,ηρ₁₁]。激发态布居指数衰减，流失的概率全部进入基态，非对角相干幅随√η衰减而非η。
\par\noindent\textbf{(3)} Bloch z=ρ₁₁−ρ₀₀按 z(t)=η(z(0)+1)−1 演化，横向 x,y按√η衰减。纵向弛豫率Γ对应T₁=1/Γ，横向相干率Γ/2对应T₂=2/Γ=2T₁；若另有纯退相干γ\_φ，则1/T₂=1/(2T₁)+γ\_φ。
\par\noindent\textbf{(4)} 初始|Φ⁺⟩=(|00⟩+|11⟩)/√2，对第二比特通道后的态含无跳跃向量(|00⟩+√η|11⟩)/√2及跳跃分量√(1−η)|10⟩/√2。与原Bell态的保真度 F=⟨Φ⁺|ρ\_out|Φ⁺⟩=(1+√η)²/4，η=1给1、η→0给1/4，核对衰减使最大纠缠损失。
